\section{Cơ sở toán học của lan truyền ngược}
\label{sec:cosotoan}

\subsection{Quy ước ký hiệu}
\label{subsec:kyhieu}

Toàn bộ tài liệu tuân thủ nghiêm ngặt quy ước sau. Mọi kết luận ở các mục sau chỉ
đúng khi quy ước này được giữ nguyên, đặc biệt là vị trí của các phép chuyển vị.

\begin{itemize}[leftmargin=*, itemsep=2pt]
  \item \textbf{Gradient là vector cột} (bố cục mẫu số, \emph{denominator
    layout}). Với $f:\mathbb{R}^n\to\mathbb{R}$, ta có
    $\partial f/\partial \mathbf{x}\in\mathbb{R}^{n\times 1}$.
  \item Mạng có $L$ tầng. Tầng $l$ có $n_l$ nơ-ron;
    $\mathbf{W}^{[l]}\in\mathbb{R}^{n_l\times n_{l-1}}$ và
    $\mathbf{b}^{[l]}\in\mathbb{R}^{n_l}$.
  \item Lan truyền tiến:
    \begin{equation}
      \mathbf{z}^{[l]} = \mathbf{W}^{[l]}\mathbf{a}^{[l-1]} + \mathbf{b}^{[l]},
      \qquad
      \mathbf{a}^{[l]} = \sigma\!\left(\mathbf{z}^{[l]}\right),
      \qquad
      \mathbf{a}^{[0]} = \mathbf{x}.
      \label{eq:forward}
    \end{equation}
    Hàm $\sigma$ tác động theo từng phần tử.
  \item Tín hiệu lỗi tại tầng $l$:
    \begin{equation}
      \boldsymbol{\delta}^{[l]}
      \;:=\; \frac{\partial\mathcal{L}}{\partial\mathbf{z}^{[l]}}
      \;\in\;\mathbb{R}^{n_l}.
      \label{eq:deltadef}
    \end{equation}
\end{itemize}

\begin{notebox}
$\boldsymbol{\delta}$ được định nghĩa theo \emph{tiền kích hoạt} $\mathbf{z}$,
không phải theo $\mathbf{a}$. Đây là lựa chọn có chủ đích: nó khiến ma trận đường
chéo chứa $\sigma'$ xuất hiện ở bên trái trong công thức đệ quy, và làm lộ rõ cấu
trúc tích mà toàn bộ tài liệu xoay quanh.
\end{notebox}

\subsection{Jacobian giữa hai tầng liên tiếp}

Từ~\eqref{eq:forward}, viết $\mathbf{z}^{[k]}$ như hàm của
$\mathbf{z}^{[k-1]}$:
\[
  \mathbf{z}^{[k]}
  = \mathbf{W}^{[k]}\sigma\!\left(\mathbf{z}^{[k-1]}\right) + \mathbf{b}^{[k]}.
\]
Lấy đạo hàm riêng từng phần tử, với $i$ chỉ số hàng của $\mathbf{z}^{[k]}$ và $j$
chỉ số hàng của $\mathbf{z}^{[k-1]}$:
\[
  \frac{\partial z^{[k]}_i}{\partial z^{[k-1]}_j}
  = W^{[k]}_{ij}\,\sigma'\!\left(z^{[k-1]}_j\right).
\]
Gom lại thành ma trận, ta được Jacobian tiến
\begin{equation}
  \frac{\partial \mathbf{z}^{[k]}}{\partial \mathbf{z}^{[k-1]}}
  = \mathbf{W}^{[k]}\operatorname{diag}\!\left(\sigma'\!\left(\mathbf{z}^{[k-1]}\right)\right)
  \;\in\;\mathbb{R}^{n_k\times n_{k-1}}.
  \label{eq:jacforward}
\end{equation}
Ma trận đường chéo nằm \emph{bên phải} vì $\sigma'$ mang chỉ số $j$ của tầng
nguồn.

\subsection{Công thức đệ quy của tín hiệu lỗi}

\begin{infobox}{Định lý --- công thức đệ quy của tín hiệu lỗi}
Với quy ước ở \hyperref[subsec:kyhieu]{Mục II.1},
\begin{equation}
  \boldsymbol{\delta}^{[l]}
  = \operatorname{diag}\!\left(\sigma'\!\left(\mathbf{z}^{[l]}\right)\right)
    \left(\mathbf{W}^{[l+1]}\right)^{\!\top}\boldsymbol{\delta}^{[l+1]} .
  \label{eq:deltarec}
\end{equation}
\end{infobox}

\noindent\textit{Chứng minh.}\;
Áp dụng quy tắc dây chuyền cho từng thành phần. Mất mát $\mathcal{L}$ phụ thuộc
$\mathbf{z}^{[l]}$ chỉ thông qua $\mathbf{z}^{[l+1]}$, nên
\[
  \delta^{[l]}_i
  = \frac{\partial\mathcal{L}}{\partial z^{[l]}_i}
  = \sum_{m} \frac{\partial\mathcal{L}}{\partial z^{[l+1]}_m}\,
             \frac{\partial z^{[l+1]}_m}{\partial z^{[l]}_i}
  = \sum_{m} \delta^{[l+1]}_m\, W^{[l+1]}_{mi}\,\sigma'\!\left(z^{[l]}_i\right),
\]
trong đó bước cuối dùng~\eqref{eq:jacforward} với $k=l+1$. Nhân tử
$\sigma'(z^{[l]}_i)$ không phụ thuộc $m$ nên đưa được ra ngoài tổng:
\[
  \delta^{[l]}_i
  = \sigma'\!\left(z^{[l]}_i\right)\sum_{m} W^{[l+1]}_{mi}\,\delta^{[l+1]}_m
  = \sigma'\!\left(z^{[l]}_i\right)
    \left[\left(\mathbf{W}^{[l+1]}\right)^{\!\top}\boldsymbol{\delta}^{[l+1]}\right]_i .
\]
Tổng $\sum_m W^{[l+1]}_{mi}\delta^{[l+1]}_m$ chạy theo chỉ số \emph{hàng} $m$ của
$\mathbf{W}^{[l+1]}$, tức là tích của \emph{chuyển vị} với
$\boldsymbol{\delta}^{[l+1]}$ --- đây chính là nơi dấu chuyển vị sinh ra. Viết
phép nhân theo từng phần tử với $\sigma'$ dưới dạng nhân ma trận đường chéo bên
trái, ta thu được~\eqref{eq:deltarec}.\hfill$\square$

Từ $\boldsymbol{\delta}^{[l]}$, gradient theo trọng số có dạng tích ngoài:
\begin{equation}
  \frac{\partial\mathcal{L}}{\partial\mathbf{W}^{[l]}}
  = \boldsymbol{\delta}^{[l]}\left(\mathbf{a}^{[l-1]}\right)^{\!\top},
  \qquad
  \frac{\partial\mathcal{L}}{\partial\mathbf{b}^{[l]}}
  = \boldsymbol{\delta}^{[l]}.
  \label{eq:gradW}
\end{equation}

\begin{examplebox}{Kiểm tra chiều --- cách tự soát lỗi chuyển vị}
Một công thức lan truyền ngược sai dấu chuyển vị hầu như luôn sai cả về chiều.
Vì vậy phép kiểm tra chiều là công cụ soát lỗi rẻ nhất. Với~\eqref{eq:deltarec}:
\[
  \underbrace{\operatorname{diag}\!\left(\sigma'(\mathbf{z}^{[l]})\right)}_{n_l\times n_l}
  \cdot
  \underbrace{\left(\mathbf{W}^{[l+1]}\right)^{\!\top}}_{n_l\times n_{l+1}}
  \cdot
  \underbrace{\boldsymbol{\delta}^{[l+1]}}_{n_{l+1}\times 1}
  \;=\;
  \underbrace{\boldsymbol{\delta}^{[l]}}_{n_l\times 1}\quad\checkmark
\]
Nếu viết nhầm thành $\mathbf{W}^{[l+1]}\boldsymbol{\delta}^{[l+1]}$ (không chuyển
vị), tích $n_{l+1}\times n_l$ nhân $n_{l+1}\times 1$ không hợp lệ trừ khi
$n_l = n_{l+1}$ --- đó là lý do lỗi này thường ẩn mình trong các mạng có mọi tầng
cùng độ rộng và chỉ lộ ra khi đổi kiến trúc. Tương tự
với~\eqref{eq:gradW}: $(n_l\times 1)(1\times n_{l-1}) = n_l\times n_{l-1}$, đúng
bằng chiều của $\mathbf{W}^{[l]}$.
\end{examplebox}

\subsection{Đối tượng trung tâm: tích Jacobian}

Khai triển~\eqref{eq:deltarec} đệ quy từ tầng $L$ về tầng $l$ cho kết quả then
chốt của toàn bộ tài liệu:
\begin{equation}
  \boxed{\;
  \boldsymbol{\delta}^{[l]}
  = \left(\prod_{k=l+1}^{L}
      \operatorname{diag}\!\left(\sigma'\!\left(\mathbf{z}^{[k-1]}\right)\right)
      \left(\mathbf{W}^{[k]}\right)^{\!\top}
    \right)\boldsymbol{\delta}^{[L]}\;}
  \label{eq:central}
\end{equation}
với quy ước các thừa số xếp theo $k$ tăng dần từ trái sang phải.

Kết hợp~\eqref{eq:central} với~\eqref{eq:gradW}, ta thấy gradient của tầng $l$ là
tích của $L-l$ thừa số ma trận. Đây là một \emph{tích}, không phải một tổng. Tổng
suy giảm tuyến tính; tích suy giảm theo cấp số nhân. Toàn bộ hiện tượng triệt
tiêu và bùng nổ gradient là hệ quả trực tiếp của sự khác biệt đó, và
Mục~\ref{sec:nguyennhan} chỉ làm một việc: ước lượng độ lớn của mỗi thừa số
trong~\eqref{eq:central}.

\begin{figure}[H]
\centering
\begin{tikzpicture}[
    >=Stealth,
    node distance=2.05cm,
    lay/.style={draw, rounded corners=2pt, minimum width=1.15cm,
                minimum height=0.85cm, thick, draw=aivnGreen,
                fill=aivnGreen!8, font=\small},
    loss/.style={draw, rounded corners=2pt, minimum width=1.15cm,
                 minimum height=0.85cm, thick, draw=brandIndigo,
                 fill=brandIndigo!10, font=\small},
    fwd/.style={->, thick, aivnGreen},
    bwd/.style={->, thick, brandIndigo}
]
\node[lay] (a0) {$\mathbf{a}^{[0]}$};
\node[lay, right of=a0] (a1) {$\mathbf{a}^{[1]}$};
\node[lay, right of=a1] (a2) {$\mathbf{a}^{[2]}$};
\node[lay, right of=a2] (a3) {$\mathbf{a}^{[3]}$};
\node[lay, right of=a3] (a4) {$\mathbf{a}^{[4]}$};
\node[loss, right of=a4] (LL) {$\mathcal{L}$};

% ---- forward arrows above ----
\foreach \s/\t/\k in {a0/a1/1, a1/a2/2, a2/a3/3, a3/a4/4}{%
  \draw[fwd] (\s.north east) to[bend left=32]
      node[above, font=\scriptsize, text=aivnGreen]{$\mathbf{W}^{[\k]}$}
      (\t.north west);
}
\draw[fwd] (a4.north east) to[bend left=32] (LL.north west);

% ---- backward arrows below, opacity fading toward the input ----
\draw[bwd, opacity=1.00] (LL.south west) to[bend left=32]
    node[below, font=\scriptsize, text=brandIndigo]{$\boldsymbol{\delta}^{[4]}$} (a4.south east);
\draw[bwd, opacity=0.70] (a4.south west) to[bend left=32]
    node[below, font=\scriptsize, text=brandIndigo]{$\boldsymbol{\delta}^{[3]}$} (a3.south east);
\draw[bwd, opacity=0.45] (a3.south west) to[bend left=32]
    node[below, font=\scriptsize, text=brandIndigo]{$\boldsymbol{\delta}^{[2]}$} (a2.south east);
\draw[bwd, opacity=0.25] (a2.south west) to[bend left=32]
    node[below, font=\scriptsize, text=brandIndigo]{$\boldsymbol{\delta}^{[1]}$} (a1.south east);
\draw[bwd, opacity=0.12] (a1.south west) to[bend left=32] (a0.south east);

\node[font=\scriptsize, text=aivnGreen, above=0.95cm of a2] {lan truyền tiến};
\node[font=\scriptsize, text=brandIndigo, below=0.95cm of a2]
     {lan truyền ngược --- độ đậm giảm dần theo biên độ gradient};
\end{tikzpicture}

\begin{figurecaptionmodern}
\capfig{Đồ thị tính toán của mạng $4$ tầng. Mũi tên phía trên là lan truyền tiến;
mũi tên phía dưới là lan truyền ngược, với độ đậm giảm dần về phía đầu vào để
biểu thị sự suy giảm biên độ của $\boldsymbol{\delta}^{[l]}$ theo~\eqref{eq:central}.\label{fig:graph}}
\end{figurecaptionmodern}
\end{figure}
